In questo progetto abbiamo analizzato  ImageNet-R e ImageNet-C per rilevare la presenza di drift di dati rispetto alla baseline offerta dal dataset ImageNet-1k. Abbiamo misurato il drift sia nelle feature visive che in quelle estratte dal modello ResNet50 in fase di classificazione, rispettivamente utilizzando la distanza di Wasserstein e Maximum Mean Discrepancy (MMD). Abbiamo inoltre misurato la degradazione delle performance del modello in presenza di drift di dati.
Dai risultati ottenuti, possiamo notare che non sembra esistere una forte correlazione tra la presenza di drift nelle feature visive delle immagini e la presenza di drift nelle feature estratte nel modello. Abbiamo potuto notare che in presenza di un covariate shift dato dalla cambio di stile di rappresentazione dell'oggetto presente nell'immagine, il modello registra un calo drastico delle sue performance, arrivando ad aver un'accuratezza par a $0.3034$ e quindi registrando un decremento del $53.9\%$ rispetto alla baseline. Invece, nel caso di covariate shift introdotto da una sporcatura delle immagini, le performance modello variano in base al tipo di sporcatura applicata: la maggior parte dei tipi osservati fanno registrare al modello dei cali drastici delle performance, mentre pochi tipi riescono a far mantenere al modello prestazioni accettabili. Si può notare che, in generale, sussiste una correlazione tra il calo delle performance e il drift rilevato dal test MMD applicato sulle feature del modello. Non sembra invece sussistere una forte correlazione tra il drift delle feature del modello e il drift rilevato nelle feature visuali. Possiamo effettuare un confronto tra le performance registrate dal modello su ImageNet-R e ImageNet-C:
\begin{itemize}
    \item Per i primi due livelli di severity, le performance registrate su ImageNet-C sono migliori rispetto a quelle registrate su ImageNet-R per tutti i tipi di sporcatura.
    \item A livello di severità 3, l'unico tipo di sporcatura che registra un'accuracy inferiore a quella registrata per ImageNet-R è glass\_blur.
    \item A livello di severità 4 hanno un accuracy inferiore motion\_blur, defocus\_blur e glass\_blur e tutti i tipi di sporcatura della classe noise.
    \item A livello 5, registrano accuracy inferiore a ImageNet-R tutti i tipi di blur, contrast, elastic\_transform, tutti i tipi di noise e snow
\end{itemize}
Possiamo quindi concludere che il covariate shift introdotto da un cambio di rappresentazione dell'oggetto è molto più impattante rispetto al covariate shift introdotto dalla sporcatura dell'immagine.
\section{Sviluppi futuri}
Elenchiamo di seguito alcuni possibili sviluppi futuri di questo progetto:
\begin{itemize}
    \item Nel paper "Do ImageNet Classifiers Generalize to ImageNet", la creazione del dataset ImageNet-V2 introduce uno \textbf{data drift temporale}, dato da un campionamento temporalmente successivo rispetto alla creazione di ImageNet \cite{recht_imagenet_2019}. Pur essendo immagini generate nello stesso periodo del dataset originale, sembra che il campionamento, nonostante esso sia stato studiato per mantenere le stesse distribuzioni rispetto a ImageNet, introduca drift e quindi il calo di performance rilevato. Sarebbe quindi interessante valutare l'effettiva presenza di drift nelle feature sia visive che del modello.
    \item In questo progetto, abbiamo considerato solo dataset che introducono covariate shift. Sarebbe interessante ripetere l'analisi su dataset che introducono altri tipi di drift.
    \item In questo progetto, abbiamo considerato un solo modello. Sarebbe interessante ripetere gli esperimenti con più modelli con diverse architetture.
\end{itemize}

